\documentclass[11pt,a4paper]{article}
\usepackage{graphicx}
\usepackage{booktabs}
\usepackage{amsmath}
\usepackage[utf8]{inputenc}
\usepackage{geometry}
\usepackage[hidelinks]{hyperref}
\geometry{margin=1in}

\title{Model-Tree Pairwise SVM for LLM-DNA Robustness\\
\large Experimental Protocol and Interim Results}
\author{}
\date{\today}

\begin{document}
\maketitle

\begin{abstract}
This report evaluates how stochastic decoding affects model identification from
LLM-DNA signatures.  It compares cosine nearest-neighbour retrieval with
pairwise linear and radial-basis-function (RBF) support vector machines (SVMs)
that predict whether two models are related. Positive pairs mix identical
models with models connected in a Hugging Face model tree, so this legacy
verification task is not a like-for-like alternative to exact identification.
Because the generation sweep is incomplete, this snapshot
uses a fixed 77-model cohort shared by every reported decoding setting, avoiding
comparisons between different successful-model subsets.
\end{abstract}

\section{Task definition}
For a model $m$ and decoding setting $s=(T,p)$, the pipeline generates responses
to a fixed prompt set and converts them into a DNA vector
$\mathbf{d}_{m,s}\in\mathbb{R}^{128}$.  We evaluate two related tasks.

\paragraph{Exact-model retrieval.}
Given a query DNA $\mathbf{d}_{m,s}$, rank the model DNAs in the fixed reference
gallery by cosine distance.  A retrieval is correct when the DNA of the same
model is ranked first.  This task produces top-1, top-3, top-5, and mean
reciprocal rank (MRR).

\paragraph{Model-relatedness classification.}
Given a query/reference pair $(m_i,m_j)$, predict
\begin{equation}
 y_{ij} =
 \begin{cases}
 1, & m_i=m_j\text{ or }m_i\text{ and }m_j\text{ belong to the same exported model-tree component},\\
 0, & \text{otherwise}.
 \end{cases}
\end{equation}
The classifier receives one scalar feature,
$x_{ij}=-d_{\mathrm{cos}}(\mathbf{d}_{m_i,s},\mathbf{d}_{m_j,\mathrm{ref}})$,
so larger values indicate greater similarity.  This formulation avoids treating
only the exact same model as positive and therefore represents known model
lineage in the labels.

\section{Experimental data}

\subsection{Model set}
The requested sweep contains 100 Hugging Face models listed in
\texttt{configs/rand\_chinese\_0p3b\_7b\_rerun.jsonl}.  Evaluation uses the
intersection of models with valid DNA files in both the unsuffixed reference
gallery and the relevant sweep run. To make settings directly comparable, the
snapshot first retains runs with at least 80 available models, then intersects
their model IDs with the reference gallery. This produces one cohort of 77
models used for every reported retrieval query, SVM training run, and SVM test
run. The exact cohort and pre-filter availability counts are stored in
\texttt{temp\_top\_p\_modeltree\_common\_models.json}; no model is selected
using test performance.

\subsection{Prompt set and DNA construction}
Every model receives the same 100 synthetic classical-Chinese prompts from
\texttt{data/rand/rand\_dataset\_chinese.json}.  The dataset generator uses seed
42, and the saved extraction metadata identifies the probe set as
\texttt{rand\_chinese\_100}.  Each response is embedded with
\texttt{all-mpnet-base-v2} (768 dimensions); the 100 response embeddings are
concatenated and reduced by a seeded Gaussian random projection to one
128-dimensional DNA signature.  The response length limit is 1024 tokens and
the extraction seed is 42.  Thus, prompts and feature construction are held
fixed while only decoding parameters change.

\subsection{Decoding settings}
The planned grid uses temperatures $T\in\{0.0,0.2,0.3,0.5,0.7\}$ and
top-$p\in\{0.8,0.9,1.0\}$.  For $T=0$, sampling is disabled; for $T>0$,
sampling is enabled with the displayed top-$p$.  Each suffix
\texttt{\_tXX\_pYY\_r1} identifies its temperature, top-$p$, and repeat.

The 80-model availability rule retains the following settings:
\[
(0.0,0.8),\ (0.2,0.8),\ (0.2,0.9),\ (0.2,1.0),\
(0.3,0.8),\ (0.3,0.9),\ (0.3,1.0),\ (0.5,0.8),\
(0.5,1.0),\ (0.7,0.9).
\]
The $(0.7,0.9)$ cell uses the earlier complete \texttt{\_temp07\_run} output.
The current $(0.0,0.9)$, $(0.0,1.0)$, $(0.5,0.9)$, $(0.7,0.8)$, and
$(0.7,1.0)$ runs have 46, 3, 17, 2, and 0 available models, respectively, and
are excluded until their reruns provide adequate coverage. At $T=0$, top-$p$
is inactive because sampling is disabled; the additional deterministic cells
will therefore be a data-completeness check rather than a decoding comparison.
The reference gallery consists of the existing unsuffixed DNA records under
\texttt{out/rand\_chinese}.  Its historical response files do not record an
explicit temperature/top-$p$ pair, so it should be described as a fixed
reference gallery rather than assigned an inferred decoding setting.

\section{Model-tree labels}
The relation exporter reads Hugging Face metadata and joins base-model links
with a streaming union--find procedure.  The resulting connected components
are stored in \texttt{out/hf\_model\_tree/model\_relations.jsonl}.  The export
summary is:
\begin{itemize}
  \item 11,812 requested models and 13,330 parsed nodes;
  \item 6,282 extracted base-model edges;
  \item 790 connected components of size at least two;
  \item 618 skipped metadata requests; and
  \item external models excluded from the exported relation groups.
\end{itemize}
Connected-component membership is transitive.  This is a practical lineage
label, not a claim that every pair in a component has identical behaviour.

\section{SVM specification}

\subsection{Classifiers used in the reported experiment}
The linear scikit-learn pipeline is
\begin{center}
\texttt{StandardScaler()} $\rightarrow$
\texttt{LinearSVC(C=1.0, class\_weight='balanced', dual='auto',}
\texttt{max\_iter=20000, random\_state=42)}.
\end{center}
The RBF-kernel probe uses
\texttt{SVC(kernel='rbf', C=1.0, gamma='scale',
class\_weight='balanced', random\_state=42)}.  Both probes were enabled in the
snapshot reported here.  Because the input has only one feature (negative
cosine distance), the linear model learns a threshold in standardized distance
space, while the RBF model can learn a nonlinear decision interval along that
same scalar axis.  This kernel probe is distinct from the exact RBF-MMD
distributional experiment.

\subsection{Class balancing}
All positive pairs are retained.  Negatives are deterministically downsampled
without replacement using seed 42 and a negative-to-positive ratio of 1:1.
The same 1:1 balancing rule is used for the held-out test examples.  The saved
summaries therefore have a 50\% positive rate in both training and testing;
\texttt{class\_weight='balanced'} remains enabled as a safeguard.  Metrics on
this balanced test sample should not be interpreted as performance under the
natural, highly imbalanced prevalence of all possible model pairs.

\section{Train/test protocol}
The SVM uses leave-one-decoding-run-out evaluation.  For each displayed
temperature/top-$p$ cell:
\begin{enumerate}
  \item that run is held out and all of its valid query/reference pairs form
        the test pool;
  \item every other run supplied to the evaluator forms the training pool;
  \item training and test pools are balanced independently as described above;
        and
  \item the SVM is fit once on the training pool and evaluated on the held-out
        setting.
\end{enumerate}
Before constructing any pairs, every retained run and the reference gallery are
restricted to the same 77 model IDs. Thus, both the training pool and held-out
test pool use the identical model cohort, although they use different decoding
runs.
Training and testing are therefore performed \emph{across different decoding
settings}, not on a random pair split within the same setting.  A training run
may share the same temperature or top-$p$ with the test run, but it cannot be
the identical run specification.

There are 10 retained decoding runs. Consequently, each fold trains on nine
runs and tests on one, a 90\%/10\% split by run. For one run, the 77 query models
are paired with the same 77 reference models, producing $77^2=5{,}929$ ordered
pairs. Of these, 99 are positive: 77 same-model diagonal pairs and 22 ordered
off-diagonal pairs connected by the exported model-tree labels. The remaining
5,830 pairs are negative. Because the cohort and relation labels are fixed,
these counts are identical in every fold and for both SVM variants.

\begin{table}[ht]
\centering
\caption{Pair accounting for every leave-one-decoding-run-out fold. ``Raw
negative'' is the number available before sampling; all positives are retained,
and an equal number of negatives is sampled with seed 42. Linear and RBF SVMs
use the same sampled train and test pairs.}
\label{tab:pair-counts}
\begin{tabular}{lrrrrr}
\toprule
Split & Runs & Positive & Raw negative & Sampled negative & Final samples \\
\midrule
Train & 9 & 891 & 52,470 & 891 & 1,782 \\
Test  & 1 & 99  & 5,830  & 99  & 198 \\
\bottomrule
\end{tabular}
\end{table}

Thus, the raw pair pool is highly imbalanced: positives are only
$99/5{,}929\approx1.67\%$ of each run's pairs. The reported classification
metrics are instead computed on the balanced 198-pair test sample, with 99
positives and 99 negatives. This distinction is important when interpreting
accuracy: the reported values measure discrimination on a controlled balanced
test set, not accuracy at the natural pair prevalence.

Model identities and model-tree components are not held out: the same model can
appear under different decoding settings in training and testing.  The current
question is therefore robustness to a decoding shift for known models, not
generalization to unseen model families.  A family-held-out protocol would
require grouping by model-tree component and is a separate experiment.

\section{Metrics}

\paragraph{Retrieval metrics.}
Top-$k$ is the fraction of queries for which the exact model appears among the
first $k$ candidates.  MRR is
$\frac{1}{N}\sum_{i=1}^{N}1/\mathrm{rank}_i$.  Macro precision, recall, and F1
are also saved for exact-model top-1 predictions.

\paragraph{Binary SVM metrics.}
We report accuracy, balanced accuracy, positive-class precision, recall and F1,
specificity, Matthews correlation coefficient (MCC), ROC-AUC, and average
precision.  Balanced accuracy, MCC, ROC-AUC, and average precision are included
because ordinary accuracy alone can be misleading for related/unrelated pair
classification.  The summary files also record TP, TN, FP, FN, the original
pair counts, and the sampled pair counts.

\paragraph{Missing settings.}
A missing point in the line plots means that the run had fewer than 80 available
models when this snapshot was generated, normally because GPU scarcity left too
many generation tasks failed or unfinished. Missing points are not zero scores,
and lines are broken rather than interpolated across them. These runs are being
rerun and will be added once they satisfy the common-cohort rule.

\section{Interim results}
Table~\ref{tab:summary} reports the snapshot values stored in
\texttt{temp\_top\_p\_modeltree\_summary.json}.  These values will be replaced
after all planned sweep cells have completed.

\begin{table}[ht]
\centering
\caption{Fair per-setting comparison on the same 77-model cohort. Top-1 and MRR
evaluate exact-model cosine retrieval. Lin. and RBF columns report balanced-test
accuracy and ROC-AUC under leave-one-run-out evaluation. Runs with fewer than
80 available models are excluded from this snapshot.}
\label{tab:summary}
\resizebox{\textwidth}{!}{%
\begin{tabular}{lrrrrrrrrr}
\toprule
Run & $T$ & Top-$p$ & Models & Top-1 & MRR & Lin. acc. & Lin. AUC & RBF acc. & RBF AUC \\
\midrule
temp00\_p08\_r1 & 0.0 & 0.8 & 77 & 0.922 & 0.935 & 0.808 & 0.906 & 0.843 & 0.837 \\
temp02\_p08\_r1 & 0.2 & 0.8 & 77 & 0.442 & 0.532 & 0.722 & 0.808 & 0.783 & 0.764 \\
temp02\_p09\_r1 & 0.2 & 0.9 & 77 & 0.455 & 0.541 & 0.753 & 0.834 & 0.773 & 0.792 \\
temp02\_p10\_r1 & 0.2 & 1.0 & 77 & 0.442 & 0.540 & 0.732 & 0.813 & 0.758 & 0.787 \\
temp03\_p08\_r1 & 0.3 & 0.8 & 77 & 0.377 & 0.481 & 0.753 & 0.774 & 0.732 & 0.799 \\
temp03\_p09\_r1 & 0.3 & 0.9 & 77 & 0.351 & 0.468 & 0.737 & 0.798 & 0.732 & 0.795 \\
temp03\_p10\_r1 & 0.3 & 1.0 & 77 & 0.260 & 0.353 & 0.687 & 0.713 & 0.682 & 0.752 \\
temp05\_p08\_r1 & 0.5 & 0.8 & 77 & 0.247 & 0.333 & 0.657 & 0.708 & 0.646 & 0.731 \\
temp05\_p10\_r1 & 0.5 & 1.0 & 77 & 0.169 & 0.255 & 0.606 & 0.633 & 0.616 & 0.638 \\
temp07\_p09\_legacy & 0.7 & 0.9 & 77 & 0.078 & 0.170 & 0.561 & 0.533 & 0.545 & 0.581 \\
\bottomrule
\end{tabular}
}
\end{table}

On the fixed cohort, exact-model retrieval falls sharply as decoding becomes
more stochastic: top-1 decreases from 0.922 at $T=0$ to 0.442--0.455 at
$T=0.2$, 0.260--0.377 at $T=0.3$, 0.169--0.247 at $T=0.5$, and 0.078 for the
legacy $T=0.7$ run. The pairwise probes degrade more gradually. This does not
establish that model-tree relatedness survives better than exact identity: 77
of the 99 positive ordered pairs are exact same-model pairs, the binary task is
much easier than 77-way retrieval, and its test prevalence is artificially
50\%. The supported statement is only that aggregate same-or-related distance
separation degrades more slowly on this controlled balanced sample.

The RBF kernel does not consistently dominate the linear threshold. Across the
fixed-cohort cells, their accuracies differ by at most 0.061: RBF is higher at
$T=0$, $T=0.2$, and $(0.5,1.0)$, while linear is slightly higher for the
$T=0.3$ cells, $(0.5,0.8)$, and the legacy $T=0.7$ cell. This mixed difference is
consistent with both classifiers receiving only one scalar distance feature;
there is limited structure for the kernel to exploit.  These observations
remain provisional because each cell has one generation repeat and the grid is
not complete.

\section{Figures}

\begin{figure}[ht]
  \centering
  \includegraphics[width=0.95\linewidth]{out/dna_probe_sweep_modeltree/temp_top_p_modeltree_curves.png}
  \caption{Exact-model retrieval on the shared 77-model cohort. Each panel fixes
  top-$p$; labeled lines show top-1, top-3, top-5, and MRR as temperature changes.
  Missing settings failed the 80-model availability rule, and lines are broken
  rather than interpolated across those settings.}
  \label{fig:retrieval-grid}
\end{figure}

\begin{figure}[ht]
  \centering
  \includegraphics[width=0.95\linewidth]{out/dna_probe_sweep_modeltree/temp_top_p_modeltree_curves_probe_metrics.png}
  \caption{Pairwise-SVM accuracy on the shared 77-model cohort. Each panel fixes
  top-$p$ and compares the labeled linear and RBF-SVM lines across temperature.
  Each point tests the displayed run after training on all other retained runs.
  Accuracy equals balanced accuracy here because the test pairs are sampled
  with equal positive and negative counts.}
  \label{fig:svm-grid}
\end{figure}

\section{Reproducibility and limitations}
The evaluation implementation is in
\texttt{scripts/evaluate\_dna\_probe\_sweep.py}; model-tree export is in
\texttt{scripts/export\_hf\_model\_tree.py}.  Machine-readable summaries and
predictions are stored under \texttt{out/dna\_probe\_sweep\_modeltree}.
Random projection, negative sampling, and SVM fitting all use seed 42.

The main limitations of the interim experiment are the incomplete decoding
grid, the provisional 80-model inclusion threshold, one repeat per included
cell, test-set downsampling, transitive model-tree labels, and model-family
overlap between training and testing. The historical reference gallery also
lacks recorded decoding parameters, so the temperature curve measures mismatch
to that particular gallery rather than a clean causal temperature effect. The
fixed cohort removes unequal coverage as a comparison confound but reduces the
population from 100 to 77 models.

This table is therefore retained as a legacy diagnostic, not the primary
experiment. The revised protocol uses a predeclared 300-model cohort, four
independent generation seeds, explicit repeat galleries, full-natural-
prevalence verification, multiclass SVMs on the full DNA vector, and bounded
unknown-temperature evaluation. Its executable specification is
\texttt{docs/revised\_experiment\_protocol.md}.

\end{document}
